\title{Parameter-Efficient Orthogonal Finetuning via Factorization}

\begin{abstract}
While large foundation models have become widespread, the computational cost of training them from scratch is substantial. As a result, methods that adapt these models to specific downstream tasks in a resource-efficient manner are increasingly sought after. In this work, we investigate a systematic finetuning approach known as Orthogonal Finetuning (OFT) for transferring foundation models to target tasks. Although OFT exhibits strong generalization, its use of high-dimensional orthogonal matrices results in a significant number of trainable parameters. To mitigate this limitation, we first analyze OFT through the lens of information transmission and identify essential properties for enhancing parameter-efficiency. Drawing inspiration from the Cooley-Tukey fast Fourier transform, which is renowned for its efficient information flow, we devise an efficient orthogonal parameterization based on butterfly structures. Incorporating this structure into OFT yields a new, parameter-efficient finetuning approach—Orthogonal Butterfly (BOFT). BOFT extends the OFT framework by providing a generalized orthogonal finetuning methodology. We evaluate our method extensively by adapting large vision transformers, language models, and text-to-image diffusion models to a variety of vision and language downstream tasks.
\end{abstract}

\section{Introduction}

Models like ChatGPT~\cite{brown2020language,chen2021evaluating} and Stable Diffusion~\cite{rombach2022high} exemplify the superior generalization capability of large-scale foundation models. This leap in performance is accompanied by a steep rise in parameter counts 

(\emph{e.g.}, GPT-3 contains approximately 175B parameters). Consequently, developing foundation models from the ground up has become an increasingly resource-intensive undertaking. The continued advancement of the field thus hinges on the ability to repurpose these powerful pre-trained models, rather than rebuilding them anew. Efficient adaptation strategies for existing foundation models to new tasks are essential. Current adaptation techniques can be broadly categorized into three types: \emph{model finetuning}~\cite{hulora2022,zaken2022bitfit,guo2020parameter,raffel2020exploring,qiu2023controlling,zhang2022adaptive,chavan2023one}, where a selected portion of the parameters of an unchanged architecture are updated; \emph{adapter tuning}~\cite{rebuffi2017learning,houlsby2019parameter,pfeiffer2020adapterfusion,lin2020exploring,he2021towards}, which integrates additional trainable components into the existing model; and \emph{prompt tuning}~\cite{li2021prefix,lester2021power}, involving the introduction of learnable prefix tokens to model inputs. Of these, model finetuning stands out for its simplicity and efficacy, while also maintaining inference speed without additional latency.

The central concept of model finetuning is to maintain a degree of similarity between the original pretrained model and its finetuned version according to specific criteria, thereby retaining the knowledge acquired during pretraining. Traditional approaches to finetuning, for example, typically rely on a reduced learning rate; this practical but heuristic method helps minimize deviations between the pretrained and updated models. Given the structured composition of weight matrices, more systematic techniques have been developed that focus on preserving the relationships within the matrices—particularly, maintaining the pairwise angular relationships between neurons. This observation has inspired the Orthogonal Finetuning~(OFT)~\cite{qiu2023controlling} framework, which guarantees that all neurons within a given layer are subjected to identical orthogonal transformations, thereby ensuring the invariance of pairwise angles among neurons during the finetuning procedure. Although OFT has achieved favorable outcomes in terms of both generalizability and convergence for text-to-image diffusion model finetuning~\cite{qiu2023controlling}, one significant drawback is the substantial number of parameters involved, as orthogonal matrices in high-dimensional spaces are inherently large. To mitigate this, OFT employs a block-diagonal structure in the orthogonal matrix, substantially decreasing the parameter count. Nonetheless, this gain in parameter efficiency introduces certain limitations: the block-diagonal orthogonal matrix enforces a rigid sparsity pattern and compartmentalizes orthogonal transformations to separate blocks. Such a design, despite its efficiency, can lead to suboptimal inductive biases (\emph{e.g.}, a reduction in expressivity and the inability to replicate classic linear transformations), and critically, there remains an open question regarding the optimal selection of the block-wise sparsity pattern.

A central challenge in this context lies in constructing a dense orthogonal matrix while maintaining a low parameter count. Although it might initially appear impractical—since a dense orthogonal matrix of dimension $d$ traditionally necessitates $\mathcal{O}(d^2)$ parameters—we propose a distinct strategy: forming a dense orthogonal matrix through the product of several factorized sparse matrices. This method builds on the observation that one can minimize the number of trainable parameters by accepting additional computational cost in exchange for lower storage. Because we express the orthogonal matrix as a series of sparse matrix multiplications, this sequence of operations must be executed repeatedly during each training iteration. To conceptualize the process of matrix factorization, we reframe the construction of a dense orthogonal matrix as a task of information transmission. Within this paradigm, assembling a dense orthogonal matrix by multiplying sparse matrices can be viewed as the movement of information across a grid-like graph. Such a perspective not only accommodates various strategies for sparse matrix factorization—limiting the parameter count yet enabling the synthesis of expressive dense matrices—but also underpins our approach.

To optimize parameter efficiency within our information transmission model, we take cues from the butterfly networks used in the Cooley-Tukey fast Fourier transform algorithm~\cite{cooley1965algorithm}, which enable highly effective communication among network nodes~\cite{klasing1994broadcasting}. Specifically, the butterfly graph underlying the Cooley-Tukey method provides a structural template for a type of sparse matrix factorization, termed butterfly factorization. Using butterfly factorization, we can realize a dense matrix of size $d$ as the product of $\mathcal{O}(\log d)$ sparse matrices, where each sparse factor contains just $\mathcal{O}(d)$ nonzero entries. When each of these sparse matrices is constrained to be orthogonal, this results in a highly efficient orthogonal matrix parameterization using only $\mathcal{O}(d\log d)$ parameters—a drastic reduction compared to conventional methods. Leveraging this efficiency, we introduce Orthogonal Butterfly~(BOFT), a new finetuning approach characterized by its parameter efficiency. BOFT generalizes the OFT method, establishing a comprehensive orthogonal finetuning framework. Notably, both block-diagonal and butterfly matrix structures share a reliance on sparsity to curtail the effective number of trainable parameters within orthogonal matrices. Comparing our methodology to LoRA~\cite{hulora2022}, which utilizes low-rank decomposition, it becomes clear that both low-rank and sparse matrices are foundational classes of structured matrices~\cite{candes2011robust} used to improve parameter efficiency.

In contrast to OFT's use of a block-diagonal architecture to balance expressivity with regularity, BOFT leverages the butterfly structure to facilitate a more continuous interpolation between full orthogonal matrices (enhancing expressivity) and the identity matrix (ensuring regularity). This allows for the identification of a more restricted subset of the orthogonal group tailored for specific downstream tasks. Given the butterfly structure's prominence in efficient linear transformations—such as the discrete Fourier transform and discrete cosine transform—we posit that employing this structure for parameter efficiency introduces a helpful inductive bias that fosters generalization and mitigates overfitting. Our rationale is rooted in the butterfly structure's capacity to readily realize numerous canonical linear transformations, a property the block-diagonal structure in OFT fundamentally lacks. The principal contributions of our work are as follows:

\begin{itemize}[leftmargin=*,nosep]
    \item We investigate parameter efficiency within orthogonal finetuning by proposing a novel information transmission framework. This approach reinterprets the construction of a parameter-efficient dense orthogonal matrix as an information flow problem over a graph with grid topology.
    \item Drawing on the butterfly structures integral to the Cooley-Tukey algorithm, we introduce Orthogonal Butterfly, a parameter-efficient orthogonal finetuning technique conceptualized under our information transmission framework.
    \item We offer several theoretical results explaining how BOFT achieves a substantial reduction in the number of trainable parameters while maintaining high expressivity and stable optimization dynamics. The matrix factorization in BOFT also enables an effective weight interpolation scheme (refer to Figure~\ref{fig:boft_interpolation}).
    \item For the first time, we extend orthogonal finetuning to a wide range of tasks beyond the scope of controllable text-to-image generation~\cite{qiu2023controlling}, thereby demonstrating its promise as a versatile finetuning strategy. We assess BOFT across diverse adaptation settings including computer vision and natural language processing tasks, where it consistently surpasses leading parameter-efficient approaches in both efficiency and generalization.
\end{itemize}

\textbf{Parameter-efficient finetuning (PEFT)}. As foundation models (\emph{e.g.}, \cite{brown2020language,radford2021learning,rombach2022high,kirillov2023segment}) continue to expand in size and capability, there is increasing focus on adapting these models to specific tasks while minimizing the number of tunable parameters~\cite{treviso2023efficient,ding2023parameter,lialin2023scaling}. Numerous PEFT techniques have been introduced~\cite{houlsby2019parameter,aghajanyan2020intrinsic,hulora2022,edalati2022krona,wang2022adamix,gheini2021cross,zaken2022bitfit,guo2020parameter,sung2021training,ansell2022composable,lester2021power,li2021prefix,vu2022spot,he2021towards,mao2021unipelt,karimi2021compacter,liu2022few,sung2022lst,chen2022parameter,jia2022visual,chen2022adaptformer,zhang2022neural,jie2023fact,lian2022scaling,luo2023towards}, with reparameterization-based approaches~\cite{aghajanyan2020intrinsic,hulora2022,edalati2022krona} being particularly relevant to our contribution. LoRA~\cite{hulora2022} achieves strong results on natural language tasks by modifying the pretrained weights through the addition of a product of two low-rank matrices. Since LoRA assigns a fixed rank to all its low-rank components, later approaches~\cite{zhang2022adaptive,valipour2022dylora,zhang2023increlora} have sought to dynamically allocate rank across layers to optimize the use of available parameters. This simple low-rank reparameterization strategy has become widespread~\cite{dettmers2023qlora,zi2023delta,chavan2023one}. Drawing from ideas in hyperspherical energy and its relationship to generalization~\cite{liu2018learning,liu2021orthogonal}, \cite{qiu2023controlling} introduce orthogonal finetuning (OFT) as an alternative for adapting text-to-image diffusion models, wherein an orthogonal transformation is learned for each layer's neurons. OFT outperforms LoRA in terms of generalization and stability, though typically with a larger number of trainable parameters. Thus, further reducing OFT's parameter requirements is valuable. Additionally, OFT's efficacy beyond text-to-image diffusion adaptation remains to be established. BOFT addresses OFT's parameterization overhead by utilizing butterfly factorization, thereby enabling broader applications of orthogonal finetuning to diverse adaptation scenarios.

\textbf{Butterfly structures}. The radix-2 Cooley-Tukey algorithm decomposes an $N$-point discrete Fourier transform into two transforms of size $\frac{N}{2}$, generating a butterfly structure that is expressible as a sequence of sparse matrices (collectively known as a butterfly matrix). Such butterfly matrices have been adopted for parameterizing orthogonal matrices to circumvent pivoting in Gaussian elimination and enhance efficiency~\cite{parker1995random}, improve recurrent neural network training stability~\cite{jing2017tunable}, and for kernel approximation~\cite{munkhoeva2018quadrature}. Techniques developed in~\cite{dao2019learning,dao2019kaleidoscope} focus on learning rapid linear transformations through butterfly parametrization, and works like~\cite{chen2022pixelated,dao2022monarch} exploit butterfly matrices to enable efficient neural network training. Furthermore, butterfly structures contribute to fast matrix-vector multiplication~\cite{michielssen1996multilevel,o2010algorithm}, efficient data-sparse matrix representations~\cite{li2015butterfly}, and enhanced communication in network transmission~\cite{klasing1994broadcasting,li2003linear,rajasingh2016transmission}. Differing from these prior directions, we specifically utilize butterfly structures to advance the parameter-efficiency of OFT.

\section{An Information Transmission View on Orthogonal Finetuning}

We begin by outlining foundational aspects of OFT. To adapt the pretrained weight matrix, OFT reformulates the new weights as a product between a trainable orthogonal matrix and the fixed pretrained weight matrix. In contrast with LoRA, which employs an additive low-rank matrix to adjust the weights, OFT introduces a multiplicative orthogonal matrix for weight modification. While LoRA achieves parameter efficiency by leveraging a low-rank design, the original OFT attains efficiency through the use of a block-diagonal orthogonal architecture~\cite{qiu2023controlling}; parameter savings are increased by reducing the size of the diagonal blocks. Figure~\ref{fig:comp_lora} provides a visual comparison of these approaches. The rationale for orthogonal transformations in fine-tuning weight matrices is to maintain neuron pairwise angular relationships~\cite{liu2018learning,liu2021orthogonal,qiu2023controlling}, thereby largely conserving semantic information acquired during pretraining. Specifically, OFT trains an orthogonal matrix $\bm{R}\in\mathbb{R}^{d\times d}$ alongside a fixed pretrained linear layer $\bm{W}^0\in\mathbb{R}^{d\times n}$, updating the standard forward operation from $\bm{z}=(\bm{W}^0)^\top\bm{x}$ to $\bm{z}=(\bm{R}\bm{W}^0)^\top\bm{x}$, where $\bm{x}\in\mathbb{R}^{d}$ represents the input and $\bm{z}\in\mathbb{R}^{n}$ the output vector. To guarantee zero initialization, OFT sets $\bm{R}$ as the identity matrix at the start. Maintaining orthogonality during fine-tuning relies on Cayley parameterization, as proposed in~\cite{qiu2023controlling,liu2021orthogonal}; that is, $\bm{R}=(\bm{I}+\bm{Q})(\bm{I}-\bm{Q})^{-1}$ with a skew-symmetric matrix $\bm{Q}$ satisfying $\bm{Q}=-\bm{Q}^\top$. For efficient parameter usage, the orthogonal matrix $\bm{R}$ is structured in block-diagonal form: $\text{diag}(\bm{R}_1,\bm{R}_2,\cdots,\bm{R}_r)$, where each block $\bm{R}_i\in\mathcal{R}^{b\times b}$ for all $i$, with $br=d$. However, this block-diagonal configuration imposes a constraint: neuron dimensions (each column in $\bm{W}^0$) are partitioned into $r$ groups, and each group is processed independently by a different orthogonal matrix. Even though empirical results suggest the effectiveness of the block-diagonal design, the arbitrary division of neuron dimensions by index is questionable, thereby motivating interest in dense orthogonal matrices. This raises an important question: \emph{Is it possible to build a dense orthogonal matrix while still preserving parameter efficiency?}

To tackle this problem, we introduce the idea of constructing a dense orthogonal matrix as a sequential product of several sparse orthogonal matrices. In order to present a coherent yet accessible approach for examining the sparsity structures in orthogonal matrix factorizations, we reinterpret the generation of a dense orthogonal matrix in OFT within the framework of an information transmission paradigm. Concretely, forming a dense matrix $\bm{R}\in\mathbb{R}^{d\times d}$ via a sequence of $m$ square matrices such that $\thickmuskip=2mu \medmuskip=2mu \bm{R}=\bm{B}_m\bm{B}_{m-1}\cdots\bm{B}_1$ can be analogized to relaying information through a network arranged as a $d\times (m+1)$ grid, as depicted in Figure~\ref{fig:info_trans}. This perspective is motivated by the fact that a dense square matrix of dimension $d$ naturally models a fully connected mapping from a set of $d$ nodes to another set of $d$ nodes. For the entry $R_{ij}$ in the matrix $\bm{R}$, a zero indicates that communication between the $j$-th and $i$-th nodes is blocked, whereas a nonzero value signifies possible transmission. Consequently, the expression of $\bm{R}$ as a product $\bm{B}_m\bm{B}_{m-1}\cdots\bm{B}_1$ can be seen as a process of successive information exchanges dictated by the sequence of graphs encoded by each $\bm{B}_i$. The progression of information starts by following the connections in $\bm{B}_1$ and concludes with those in $\bm{B}_n$. 

As illustrated specifically in Figure~\ref{fig:info_trans}, we analyze the case where $\bm{R}=\bm{B}_5\bm{B}_4\bm{B}_3\bm{B}_2\bm{B}_1$, visualizing both the sparsity configurations and the resulting induced graph. The depicted graph corresponds to the unrolled computation of the matrix product. In this induced structure, each matrix $\bm{B}_i$ is interpreted as defining connections from nodes at the $i$-th level to those at the $(i+1)$-th level of the network. Specifically, a nonzero entry at position $(j_1, j_2)$ in $\bm{B}_i$ marks the existence of a directed link from the $j_2$-th node on level $i$ to the $j_1$-th node on level $(i+1)$ (a zero implies no such link). For the overall product $\bm{B}_5\bm{B}_4\bm{B}_3\bm{B}_2\bm{B}_1$ to yield a dense matrix, it is required that every node at the initial layer can reach all nodes at the sixth layer through a path of connections. If we restrict our focus to $\bm{R}=\bm{B}_4\bm{B}_3\bm{B}_2\bm{B}_1$, which joins the nodes from the first to the fifth layer, it becomes apparent, for example, that the first node’s information fails to reach the third node—demonstrated by the zero at position $(3,1)$ in the resulting product’s sparsity pattern. Analogous observations can be made for $\bm{R}=\bm{B}_3\bm{B}_2\bm{B}_1$ and $\bm{R}=\bm{B}_2\bm{B}_1$, where the connectivity structure directly reflects the matrix product's sparsity. In general, for $\bm{R}\in\mathbb{R}^{d\times d}$ to be fully dense, the collection of factor matrices $\bm{B}_m, \dots, \bm{B}_1$ must together specify a network of directed edges spanning a $d\times (m+1)$ grid. Each directed edge exclusively links nodes between consecutive layers (i.e., adjacent columns), ensuring that every node at the starting level is able to propagate information to every node at the final level $(m+1)$.

The information transmission perspective offers valuable insights into the sparse structures inherent in the factorization matrices used within OFT. In Figure~\ref{fig:blk_diag}, we illustrate how $\bm{R}$ in the standard OFT manifests as a block-diagonal matrix. While this configuration decreases the quantity of learnable parameters, it simultaneously precludes the formation of a fully dense matrix $\bm{R}$. Our objective is to generate a dense orthogonal matrix by composing $m$ sparse orthogonal matrices, optimizing for a minimal number of effective trainable parameters. Interpreting this within the information transmission framework, two fundamental requirements emerge: \emph{(i)} \textbf{dense connectivity}, ensuring that each initial node is linked—through at least one pathway—to every terminal node, and \emph{(ii)} \textbf{minimum free edges}, seeking to restrict the total edge count as much as possible, given the orthogonality constraints. It is important to recognize that orthogonality imposes subtle constraints on permissible edge connections between adjacent layers. Specifically, maintaining the full rank of each matrix $\bm{B}_i$ (essential for orthogonality) necessitates $d$ edges establishing a bijective mapping between nodes of the $i$-th and $(i=1)$-th levels, thereby enforcing at least $d$ inter-level edges (\emph{e.g.}, as demonstrated in Figure~\ref{fig:info_trans} with $d=4$). These particular $d$ edges are required by orthogonality and are excluded from the trainable edge count, since the corresponding entries are not subject to optimization (\emph{e.g.}, in a $d\times d$ orthogonal matrix with $d$ nonzero values, those entries must be fixed as $\pm 1$). The reduction in parameter count inherent to orthogonal matrices introduces additional interdependence between edges: for instance, setting a single zero entry (such as at the $(1,1)$ location in a $2\times 2$ orthogonal matrix) may force another zero (e.g., at $(2,2)$), meaning that eliminating one edge can imply the removal of another. Consequently, the set of feasible edge configurations can be altered by orthogonality constraints, either requiring or permitting the addition or removal of certain edges. Thus, examining the nonzero pattern in the constructed orthogonal matrix through the lens of the information transmission model proves especially valuable in OFT settings, as only the arrangement of nonzero trainable elements in $\bm{R}$ is consequential, whereas their specific magnitudes are irrelevant.

A straightforward fully-connected setup between two layers involves $\mathcal{O}(d^2)$ edges (that is, it requires a single dense orthogonal matrix), specifically $d^2-d$ edges (for $d=4$, this results in 12 edges). Figure~\ref{fig:info_trans} presents an instance of a possible matrix factorization, which only utilizes $10$ edges in total—fewer than what is needed for a fully dense orthogonal matrix. This approach allows us to examine OFT’s parameter-efficiency through a graphical lens and facilitates the generation of admissible factorizations within this framework. We draw on an intriguing network topology from the Cooley-Tukey algorithm, namely butterfly graphs~\cite{cooley1965algorithm}, which offer a way to connect $d$ sources and $d$ destinations densely but with just $\mathcal{O}(d\log d)$ edges. To illustrate, the layout in Figure~\ref{fig:info_trans} achieves dense communication using $10$ edges, but a butterfly network reduces this to only $8$. The subsequent section outlines how utilizing the butterfly configuration further enhances parameter efficiency.

\section{Orthogonal Parameterization by Butterfly Factorization}

The butterfly architecture has its origins in the Cooley-Tukey algorithm for executing the fast Fourier transform. In the context of Fourier transforms, a modification in the frequency space can propagate a widespread impact in the spatial domain. This characteristic aligns with the challenges in our information transmission setting—any node in the input layer is able to communicate with all nodes at the output layer. The butterfly topology has also been adopted in computer network design~\cite{solihin2015fundamentals,li2011linear} to promote efficient data transfer. For $k \geq 2$ and $k$ a power of two, we define the \emph{butterfly factor} $\bm{B}^F(k)$ as follows:

\begin{equation}\label{eq:but_factor}
\begin{aligned}
    \bm{B}^F(k)=\begin{bmatrix}
    \text{diag}(\bm{d}_1) & \text{diag}(\bm{d}_2) \\
    \text{diag}(\bm{d}_3) & \text{diag}(\bm{d}_4)
    \end{bmatrix}\in\mathbb{R}^{k\times k},
\end{aligned}
\end{equation}

where each $\bm{d}_i\in\mathbb{R}^{\frac{k}{2}}$ for $i = 1, \ldots, 4$ are vectors. Next, let $d=2^N$; the $d$-dimensional \emph{butterfly matrix} $\bm{B}(d)\in\mathbb{R}^{d\times d}$ is then constructed recursively by

\begin{equation}
\begin{aligned}
    \!\!\bm{B}(d)=\tilde{\bm{B}}(d,d)\!\cdot\!\begin{bmatrix}
    \bm{B}_1(\frac{d}{2})\!\!\!\!\! & \bm{0} \\
    \bm{0}\!\!\!\!\! &\bm{B}_2(\frac{d}{2})
    \end{bmatrix}= \tilde{\bm{B}}(d,d)\tilde{\bm{B}}(d,\frac{d}{2})\cdots\tilde{\bm{B}}(d,2),
\end{aligned}
\end{equation}

Here, $\bm{B}_1(\frac{d}{2})$ and $\bm{B}_2(\frac{d}{2})$ denote two butterfly matrices each of dimension $\frac{d}{2}$. We introduce the notion of a \emph{butterfly component}, expressed as $\thickmuskip=2mu \medmuskip=2mu \tilde{\bm{B}}(d,k)=\text{diag}(\bm{B}^F_1(k),\cdots,\bm{B}^F_{d/k}(k))$, which forms a $d\times d$ block-diagonal matrix with each diagonal block having size $k$. Here, every $\bm{B}^F_i(k)$ is a butterfly factor as presented in Equation~\ref{eq:but_factor}. This setup allows us to parameterize an orthogonal matrix through the butterfly matrix structure. The essential requirement is that each factor $\tilde{\bm{B}}(d,k)$ composing $\bm{B}(d)$ maintains orthogonality. Focusing on the specific instance $\tilde{\bm{B}}(d,2)$—the block-diagonal case with block size $2$—one can readily impose orthogonality via the Cayley transform or via $2$-dimensional rotations to represent each block, analogous to the approach taken in \cite{liu2021orthogonal,qiu2023controlling}. Moreover, as the sparsity pattern of every butterfly component is simply a permutation of that found in $\tilde{\bm{B}}(d,2)$, all such components are equally amenable to orthogonal parameterization. As a result, orthogonal matrices can be constructed efficiently from multiple $2\times 2$ orthogonal submatrices. Following the extension in \cite{chen2022pixelated}, we generalize the construction by defining a block butterfly component, $\tilde{\bm{B}}^b(d,k)$, where each $\bm{d}_i$ is substituted by a $b\times b$ matrix. To ensure that $\tilde{\bm{B}}^b(d,2)$ remains orthogonal, every $2b\times 2b$ block is required to be orthogonal. The support structure of any $\tilde{\bm{B}}^b(d,k)$ for $k>2$ arises by permuting the block support of $\tilde{\bm{B}}^b(d,2)$, thereby admitting the same orthogonal construction. Collectively, this forms the basis for the forward pass in BOFT, which is given as follows:

\begin{equation}
    \begin{aligned}\nonumber
    \bm{z}=\big(\bm{R}(m,b)\cdot\bm{W}^0\big)^\top\bm{x},~~~~\text{s.t.}~\bigg{\{}\bm{R}(m,b)=\prod_{i=1}^m \tilde{\bm{B}}^b_{(d,i)}~~\&~~\underbrace{\big(\tilde{\bm{B}}^b_{(d,j)}\big)^\top\tilde{\bm{B}}^b_{(d,j)}=\tilde{\bm{B}}^b_{(d,j)}\big(\tilde{\bm{B}}^b_{(d,j)}\big)^\top\!=\bm{I}_d}_{\forall j\in [1, m]}\bigg{\}},
    \end{aligned}
\end{equation}

In the above, for clarity, we abbreviate $\tilde{\bm{B}}^b(d,2^{m - i+1})$ as $\tilde{\bm{B}}^b_{(d,i)}$. Here, $\bm{I}_d$ indicates the $d$-dimensional identity matrix. The orthogonal transformation $\bm{R}(m,b)\in\mathbb{R}^{d\times d}$ is defined as a product of several orthogonal butterfly matrices. For brevity, we use the notation BOFT($m$,$b$) to refer to the method with $\bm{R}(m,\frac{b}{2})$, provided that $b\geq 2$. When $\thickmuskip=2mu \medmuskip=2mu m=1$, BOFT($1$,$b$) specializes to the block-diagonal OFT~\cite{qiu2023controlling} with block size $b$. If $b=d$, BOFT($1$,$d$) collapses to the unconstrained OFT~\cite{qiu2023controlling}, utilizing a full orthogonal matrix. Additionally, BOFT($\log \frac{2d}{b}$,$b$) employs the block butterfly structure $\bm{B}^{\frac{b}{2}}(d)$ as $\bm{R}$, resulting in a dense orthogonal mapping. Generally, the BOFT($m$,$b$) approach introduces $\frac{1}{2}(b-1)dm$ learnable parameters for fine-tuning a linear mapping of size $d\times n$. Selecting a butterfly configuration ($m=\log d, b=2$), BOFT requires only $\mathcal{O}(d\log d)$ parameters. By contrast, a dense, full orthogonal matrix in the original OFT needs $\mathcal{O}(d^2)$ parameters, and a block-diagonal OFT with $r$ blocks uses $\mathcal{O}(bd)$. Consequently, while original OFT necessitates $b=d$ to produce a dense orthogonal transformation, BOFT achieves this for arbitrary $b$.

\textbf{Identity Initialization in BOFT}. Fine-tuning techniques typically begin from the exact weights of a pretrained model to ensure the updated model does not stray far from the original. For instance, LoRA sets the initial low-rank matrices to zero. In BOFT, all butterfly components are initialized as identity matrices (i.e., using zero-initialization in the skew-symmetric matrix within the Cayley transform).

\textbf{Multiplicative Dropout in BOFT}. In LoRA~\cite{hulora2022}, Dropout is applied to the low-rank updates to mitigate overfitting, leveraging standard Dropout~\cite{srivastava2014dropout}, which is compatible with additive modifications like LoRA’s but unsuitable for the multiplicative structure in BOFT. To overcome this, we design a multiplicative Dropout for BOFT: since $\bm{R}(m,b)$ comprises $m$ orthogonal butterfly matrices that can be permuted into $2b\times 2b$ block-diagonal orthogonal structures, multiplicative Dropout randomly selects $p_1$ fraction of butterfly components and, within each, $p_2$ fraction of the diagonal blocks, then substitutes the chosen blocks by identity matrices.

\section{Intriguing Insights and Discussions}

\textbf{Expressivity of BOFT}. The butterfly structure, together with permutations, is capable of exactly representing various well-known fast linear transforms~\cite{dao2019learning,dao2019kaleidoscope} (\emph{e.g.}, fast Fourier transform, Hadamard transform). However, the degree to which our orthogonal butterfly matrix can approximate a general orthogonal matrix is not fully understood. To investigate this, we carried out simulations aiming to approximate a random dense orthogonal matrix~\cite{anderson1987generation} of dimension $512\times 512$. The data shown in Figure~\ref{fig:para_eff} represent averages over 10 independent random initializations. On the y-axis, the plot reports the approximation error, while the x-axis indicates the number of effective trainable parameters. Each colored curve corresponds to a BOFT configuration with fixed block size, and the leftmost points illustrate the error associated with BOFT($1$,$b$) (\emph{i.e.}, the standard block-diagonal OFT with block size $b$). Empirically, BOFT achieves greater parameter efficiency compared to OFT. For instance, BOFT($9$,$2$) exhibits higher expressiveness than BOFT($1$,$16$), yet requires significantly fewer parameters. Typically, configurations with smaller $b$ and larger $m$ in BOFT tend to be more parameter efficient. As an example, BOFT($6$,$4$) achieves comparable approximation accuracy to BOFT($2$,$16$) with a substantially reduced parameter count. Generally, the butterfly matrix forms a more structured subset within the orthogonal group than the block-diagonal alternative, providing theoretical grounds for the superior expressivity of BOFT over OFT for a given block size.

\begin{theorem}[Expressivity of BOFT]\label{thm:exp_boft}
    BOFT possesses greater expressive power than OFT when comparing identical block sizes. To approximate all orthogonal matrices of dimension $d$ using butterfly matrices, one can construct products of the form $\bm{B}_{d-1,1}(d)\bm{B}^\top_{d-1,2}(d)\cdots\bm{B}_{1,1}(d)\bm{B}^\top_{1,2}(d)$, where each $\bm{B}_{i,j}(d),\forall i,\forall j$ denotes a butterfly matrix.
\end{theorem}

Theorem~\ref{thm:exp_boft} highlights a straightforward path to generalize BOFT—specifically, the ultimate orthogonal matrix can be generalized as $\thickmuskip=2mu \medmuskip=2mu \bm{R}^G(m_1,b_1,m_2,b_2,l)=\bm{R}_{l,1}(m_1,b_1)\bm{R}_{l,2}^T(m_2,b_2)\cdots\bm{R}_{1,1}(m_1,b_1)\bm{R}_{1,2}^T(m_2,b_2)$, with $\bm{R}_{i,j}^T(m,b)$ being the orthogonal matrices deployed within BOFT. In the special case where $m_1=m_2=\log d$, $b_1=b_2=2$, and $l=d-1$, $\bm{R}^G(m_1,b_1,m_2,b_2,l)$ becomes sufficient to parameterize the full orthogonal group. This layered matrix construction is also referred to as a kaleidoscope hierarchy~\cite{dao2019kaleidoscope}. Nevertheless, it is important to acknowledge that greater expressive capacity does not always guarantee improved performance in finetuning. In fact, full finetuning—despite its universal expressivity—often produces subpar results. The core challenge in model finetuning lies in balancing expressivity against regularity. By integrating structural priors, BOFT expands the parameter space available for finetuning, which in turn enables the discovery of more favorable trade-offs between expressivity and regularity.

\textbf{Spectral properties}. Orthogonal finetuning consistently achieves superior spectral characteristics compared to LoRA, owing to its ability to exactly maintain the spectral norm of the pretrained weight matrix $\bm{W}^0$. This result can be elucidated via singular value decomposition, where $\bm{W}^0=\bm{U}\bm{\Sigma}\bm{V}^\top$ and the matrices $\bm{U}$ and $\bm{V}$ are orthogonal while $\bm{\Sigma}$ contains the singular values along its diagonal. In the case of both OFT and BOFT, the updated weights take the form $\bm{R}\bm{U}\bm{\Sigma}\bm{V}^\top$ by pre-multiplying with an orthogonal matrix $\bm{R}$, thereby leaving the largest singular value unchanged (i.e., preserving the spectral norm of $\bm{W}^0$). This property has been empirically demonstrated to significantly enhance both stability during training and model generalization~\cite{miyato2018spectral,yoshida2017spectral}. Further notable mathematical aspects are discussed in Appendix~\ref{app:sn_property}.

\textbf{Orthogonal finetuning as learning bilinear similarity}. The formulation of BOFT can be interpreted as optimizing a bilinear similarity function $\bm{w}^0_i\bm{R}\bm{x}$, where $\bm{w}^0_i$ denotes the $i$-th column (or neuron) of the pretrained weight matrix $\bm{W}^0$. From this perspective, BOFT is akin to learning a similarity matrix $\bm{R}$ that is subject to stringent regularization, specifically orthogonality, establishing a natural connection with established frameworks in distance metric learning~\cite{xing2002distance} and bilinear forms~\cite{roman2005advanced}.

\textbf{Inductive bias and generalization in BOFT}. In BOFT, the matrix $\bm{R}(m,b)$ frequently corresponds to a structured, constrained subset of the orthogonal group, which serves to restrict the hypothesis space and thereby introduces an inductive bias. We posit that this structured inductive bias, introduced via butterfly factorization, is advantageous for generalization, since it embodies shared structural patterns characteristic of well-known linear transforms—such as the discrete Fourier, discrete sine/cosine, and Hadamard transforms~\cite{dao2019kaleidoscope}. Additionally, the inherent sparsity in the matrix factorization process within BOFT potentially provides further implicit inductive biases~\cite{gunasekar2017implicit,li2018algorithmic,liu2019neural}.

\textbf{Comparison to butterfly-based sparse training}. A substantial body of research~\cite{dao2019learning,dao2019kaleidoscope,chen2022pixelated,dao2022monarch} has investigated the use of butterfly parameterization for sparse training. These prior works generally involve directly representing weight matrices using butterfly structures and training neural networks from the beginning. In particular, \cite{dao2022monarch} explores finetuning by projecting pretrained weights onto modified butterfly matrices before updating the projected weights for specific downstream applications. In contrast, BOFT adopts a fundamentally distinct finetuning paradigm by transforming model weights using weight matrices that are shared across layers.

\input{rebuttal-results}

\section{Concluding Remarks and Limitations}

In this work, we introduce Orthogonal Butterfly, a general-purpose, parameter-efficient finetuning approach for foundation models that leverages the butterfly matrix structure. The principal innovation enhancing parameter efficiency is the representation of a dense orthogonal matrix as the product of several sparse orthogonal matrices. To facilitate tractable matrix decompositions, we develop a framework inspired by graphical information transmission. Using this perspective, we establish that the butterfly pattern serves our goal of factoring dense orthogonal matrices into sparse orthogonal components. Our empirical evaluations showcase that BOFT performs robustly for finetuning tasks on large language models, substantial vision models, and text-to-image generative models, demonstrating BOFT’s versatility and efficacy.

Nevertheless, BOFT is not without shortcomings. Since its formulation involves expressing the final orthogonal matrix as a sequence of orthogonal matrix multiplications, the computational cost during training is somewhat higher compared to OFT. Addressing the challenge of reducing BOFT’s training time remains unresolved. On the upside, after finetuning, the orthogonal matrices obtained by BOFT can be integrated directly into the pretrained model, resulting in no additional inference cost. Furthermore, it is still an open question whether the butterfly network structure constitutes the optimal mechanism for information transfer. Our information transmission framework further offers the opportunity to incorporate insights from computer networking, where the effectiveness of various network topologies for information flow has been thoroughly analyzed. This connection suggests the potential to discover alternative, more efficient network architectures for composing dense orthogonal matrices.

\end{document}